% ch26-identity-tokens-sso
\section{الهُوية (\en{Identity})، التوكنات، و\en{SSO} بعمق}

في فصل التوثيق السابق عرفنا فكرة الرمزين على مستوى عالٍ: رمز وصول (\en{access
JWT}) قصير العمر يُرفَق مع كل طلب، ورمز تحديث (\en{refresh token}) أطول عمرًا
وظيفته استبدال الأول عند انتهائه. في هذا الفصل ننزل طبقةً أعمق ونجيب على أسئلة
عملية دقيقة: أين تُخزَّن التوكنات وبأيّ شكل؟ كيف نُبطلها في منتصف الجلسة؟ لماذا يرث
سياق قاعدة البيانات صنف \en{Identity} تحديدًا؟ وكيف يعمل الدخول عبر \en{Google}
و\en{Microsoft} خطوةً بخطوة؟ كل ما هنا مقتبس من الكود الحيّ في \code{server/src}،
بمُعرِّفات حقيقية.

\begin{idea}[الفكرة: هوية واحدة، مصدرها نحن]
\en{ScholarPath} يقبل عدّة طرق للدخول (كلمة مرور، \en{Google}، \en{Microsoft})،
لكنه في النهاية يحوّلها جميعًا إلى هوية محلية واحدة يملكها هو ويصدر لها توكناته
الخاصّة. أيًّا كان باب الدخول، بعد لحظة إثبات الهوية يصبح النظام هو المتحكّم الكامل
في الجلسة: من يدخل، وكم يبقى، ومتى يُطرَد.
\end{idea}

% ====================================================================
\subsection{ليه \code{ApplicationDbContext} يرث \code{IdentityDbContext} لا \code{DbContext}؟}

هذا سؤال يتكرّر في المناقشات. الجواب المختصر: لأن \en{ASP.NET Core Identity}
يحتاج مجموعة جداول جاهزة (المستخدمون، الأدوار، الـ\en{claims}، تسجيلات الدخول
الخارجية، التوكنات، وربط المستخدم بالدور)، والوراثة من \code{IdentityDbContext}
تُعرِّف هذه الجداول وتربط علاقاتها تلقائيًّا. هذا هو سطر إعلان الصنف الحقيقيّ:

\begin{english}
\begin{lstlisting}[language={[Sharp]C}]
public class ApplicationDbContext(
    DbContextOptions<ApplicationDbContext> options,
    IMediator? mediator = null,
    IFieldEncryptionService? encryption = null)
    : IdentityDbContext<ApplicationUser, ApplicationRole, Guid>(options),
      IApplicationDbContext
{
    public DbSet<RefreshToken> RefreshTokens => Set<RefreshToken>();
    // ... business DbSets: Scholarships, Applications, Reviews, ...
}
\end{lstlisting}
\end{english}

نقرأه سطرًا سطرًا:
\begin{itemize}
\item \code{: IdentityDbContext<ApplicationUser, ApplicationRole, Guid>} — هذه
  الوراثة تُخبر \en{EF Core} بأن يبني جداول الهوية بأنواعنا نحن: المستخدم
  \code{ApplicationUser}، والدور \code{ApplicationRole}، ومفتاح كلاهما من نوع
  \code{Guid}.
\item \code{IApplicationDbContext} — منفذ طبقة \en{Application}؛ يجعل المُعالِجات
  ترى السياق كواجهة مجرّدة لا كصنف ملموس (اتّساقًا مع \en{Clean Architecture}).
\item \code{DbSet<RefreshToken>} وبقيّة الجداول — هي جداول العمل التي نضيفها نحن
  فوق جداول الهوية.
\end{itemize}

الفرق الجوهريّ يظهر داخل \code{OnModelCreating}. صنفنا يستدعي أولًا
\code{base.OnModelCreating(builder)} فتُبنى جداول الهوية السبعة بعلاقاتها، ثم نضيف
تهيئتنا فوقها:

\begin{english}
\begin{lstlisting}[language={[Sharp]C}]
protected override void OnModelCreating(ModelBuilder builder)
{
    builder.ApplyConfigurationsFromAssembly(Assembly.GetExecutingAssembly());
    base.OnModelCreating(builder);   // <-- builds ALL Identity tables + keys
    // ... our own model tweaks after Identity is wired
}
\end{lstlisting}
\end{english}

\begin{why}
لو ورثنا \code{DbContext} العاديّ لاضطُررنا لتعريف كل جدول هوية يدويًّا: جدول
المستخدمين، وجدول الأدوار، وجدول ربط المستخدم بالدور، وجداول الـ\en{claims}
وتسجيلات الدخول الخارجية والتوكنات — مع مفاتيحها الأجنبية وفهارسها الفريدة (مثل
البريد المُطبَّع الفريد). هذا عمل طويل وسهل الخطأ. الوراثة من
\code{IdentityDbContext} تعطينا كل ذلك مضبوطًا ومختبَرًا من مايكروسوفت، ونكتفي نحن
بإضافة جداول العمل.
\end{why}

الجدول التالي يلخّص جداول الهوية التي تأتي مجّانًا مع الوراثة:

\begin{center}
\begin{tabular}{L{4.2cm} L{10.3cm}}
\toprule
{\bfseries\color{spPrimary}الجدول} & {\bfseries\color{spPrimary}ما يخزّنه} \\
\midrule
\en{AspNetUsers} & الحسابات: البريد، كلمة المرور المُجزَّأة، عدّاد المحاولات، القفل. \\
\en{AspNetRoles} & الأدوار المتاحة في المنصّة. \\
\en{AspNetUserRoles} & ربط كل مستخدم بأدواره (علاقة كثير-لكثير). \\
\en{AspNetUserClaims} & \en{claims} إضافية مرتبطة بالمستخدم. \\
\en{AspNetUserLogins} & تسجيلات الدخول الخارجية (\en{Google}/\en{Microsoft}) المربوطة بالحساب. \\
\en{AspNetUserTokens} & توكنات تُخزّنها \en{Identity} لأغراض مثل إعادة التعيين. \\
\bottomrule
\end{tabular}
\end{center}

% ====================================================================
\subsection{نموذج الأدوار والـ\en{claims}: إزاي أهندل الرولز}

\en{ScholarPath} يعرّف خمسة أدوار: \code{SuperAdmin} و\code{Admin}
و\code{Student} و\code{Consultant} و\code{ScholarshipProvider} (مزوّد المنح،
وكان اسمه سابقًا \en{Company}). المهم أن نفهم كيف تنتقل هذه الأدوار من قاعدة
البيانات إلى داخل الطلب.

\begin{keyp}[نقطة مهمة: الدور يُحمَل داخل التوكن]
لا يرجع النظام إلى قاعدة البيانات في كل طلب ليسأل «ما أدوار هذا المستخدم؟». بدلًا
من ذلك، لحظة الإصدار تُوضَع الأدوار داخل الـ\en{JWT} على هيئة \en{claims} موقّعة.
أي طلب لاحق يحمل التوكن يحمل معه أدواره، فيصبح فحص الصلاحية سريعًا وبلا استعلام.
\end{keyp}

هذا مقتطف حقيقيّ من \code{TokenService.CreateAccessToken} يبني قائمة الـ\en{claims}:

\begin{english}
\begin{lstlisting}[language={[Sharp]C}]
var claims = new List<Claim>
{
    new(JwtRegisteredClaimNames.Sub,   user.Id.ToString()),
    new(JwtRegisteredClaimNames.Email, user.Email ?? string.Empty),
    new("account_status", user.AccountStatus.ToString()),
    new(ClaimTypes.Name, user.FullName),
};
if (!string.IsNullOrEmpty(activeRole))
    claims.Add(new Claim("active_role", activeRole));   // the session's ACTIVE role
foreach (var r in roles)
    claims.Add(new Claim(ClaimTypes.Role, r));          // every role the user holds
\end{lstlisting}
\end{english}

الشرح:
\begin{itemize}
\item \code{Sub} و\code{Email} و\code{Name} — هوية المستخدم الأساسية داخل التوكن.
\item \code{active\_role} — نقطة ذكية: المستخدم قد يملك أكثر من دور (طالب ومستشار
  مثلًا)، لكن جلسته الحالية تعمل بدور واحد فعّال. نضع هذا الدور في \en{claim}
  مستقلّ.
\item \code{foreach ... ClaimTypes.Role} — نضيف كل أدوار المستخدم كـ\en{claims}
  من نوع الدور، فتُعرَف صلاحياته كاملة عند الحاجة.
\end{itemize}

الآن السؤال العمليّ: كيف تُحمى نقطة النهاية بهذا؟ نُهيّئ التحقّق ليعتبر \en{claim}
اسمه \code{active\_role} هو «claim الدور»، ثم نضع سمة \code{[Authorize(Roles=...)]}
فوق المتحكّم:

\begin{english}
\begin{lstlisting}[language={[Sharp]C}]
opts.TokenValidationParameters = new TokenValidationParameters
{
    ValidAlgorithms = [SecurityAlgorithms.RsaSha256],
    RoleClaimType   = "active_role",   // [Authorize(Roles=...)] checks the ACTIVE role
};

[Authorize(Roles = "Admin,SuperAdmin")]   // e.g. AdminController
public sealed class AdminController : ControllerBase { /* ... */ }
\end{lstlisting}
\end{english}

الشرح:
\begin{itemize}
\item \code{RoleClaimType = "active\_role"} — يجعل \code{[Authorize(Roles=...)]}
  يفحص الدور الفعّال للجلسة لا مجموعة الأدوار كلها. فمستخدم يملك دورَي طالب ومستشار،
  ويعمل الآن كطالب، لا يستطيع فتح صفحات المستشار حتى يبدّل دوره.
\item \code{[Authorize(Roles = "Admin,SuperAdmin")]} — هذه هي الحماية الفعلية على
  \code{AdminController} في الكود: فقط من دوره الفعّال \en{Admin} أو
  \en{SuperAdmin} يمرّ.
\end{itemize}

الرسم التالي يوضّح كيف «يتفرّع» التوكن الواحد إلى عدّة \en{claims}، وأيّها يستعمله
التفويض:

\begin{center}\begin{LTR}
\begin{tikzpicture}[node distance=0.55cm]
  \node[pbox, text width=2.7cm, minimum height=1.1cm] (jwt) {Access JWT\\(RS256, 15 min)};
  \node[sbox, text width=3.2cm] (c1) at (5.4,2.3)  {sub = user id};
  \node[sbox, text width=3.2cm] (c2) at (5.4,1.35) {email};
  \node[sbox, text width=3.2cm] (c3) at (5.4,0.4)  {account status};
  \node[tbox, text width=3.2cm] (c4) at (5.4,-0.55){active-role};
  \node[sbox, text width=3.2cm] (c5) at (5.4,-1.5) {role: Student};
  \node[sbox, text width=3.2cm] (c6) at (5.4,-2.45){role: Consultant};
  \foreach \c in {c1,c2,c3,c4,c5,c6} \draw[arr] (jwt.east) -- (\c.west);
  \node[fbox, text width=3.0cm, right=0.5cm of c4] (use) {used by\\Authorize(Roles)};
  \draw[arrd] (c4.east) -- (use.west);
\end{tikzpicture}
\end{LTR}\end{center}

% ====================================================================
\subsection{تخزين التوكنات: ليه الـ\en{refresh} مُجزَّأ بـ\en{SHA-256}؟}

رمز الوصول لا يُخزَّن أصلًا في القاعدة — هو موقّع ويُفحَص بمفتاح، فلا حاجة لحفظه.
أمّا رمز التحديث فيجب أن نتذكّره لنستطيع إبطاله لاحقًا. والنقطة الحسّاسة: لا نُخزّنه
كما هو، بل نُخزّن تجزئته فقط في الحقل \code{RefreshToken.TokenHash}. هذا مقتطف
حقيقيّ من \code{TokenService.IssueTokens} يوضّح لحظة الحفظ:

\begin{english}
\begin{lstlisting}[language={[Sharp]C}]
var refreshToken = GenerateSecureRefreshToken();     // 64 crypto-random bytes
var refreshHash  = HashRefreshToken(refreshToken);   // SHA-256(token) -> hex
db.RefreshTokens.Add(new RefreshToken
{
    UserId    = user.Id,
    TokenHash = refreshHash,        // the HASH is persisted, never the raw token
    ExpiresAt = refreshExpires,
    CreatedAt = now,
});
db.SaveChanges();
return new TokenPair(accessToken, refreshToken, accessExpires, refreshExpires);
\end{lstlisting}
\end{english}

الشرح:
\begin{itemize}
\item \code{GenerateSecureRefreshToken()} — يولّد 64 بايت عشوائيًّا تشفيريًّا، فالرمز
  غير قابل للتخمين إطلاقًا.
\item \code{TokenHash = refreshHash} — ما يدخل القاعدة هو تجزئة \en{SHA-256} فقط.
\item \code{return new TokenPair(..., refreshToken, ...)} — الرمز الخام يُعاد للعميل
  مرّة واحدة فقط في الاستجابة، ولا يُكتَب في القاعدة أبدًا.
\end{itemize}

عند التجديد، يأخذ النظام الرمز الخام القادم من العميل، يجزّئه بنفس الدالة، ثم يبحث
عن تجزئة مطابقة. المقارنة إذن بين تجزئتين، لا بين نصّين.

\begin{secnote}
لماذا التجزئة؟ لنفس سبب تجزئة كلمات المرور: لو تسرّبت قاعدة البيانات كاملةً، لن يجد
المهاجم إلا \en{SHA-256} لا يمكن عكسها إلى الرمز الأصلي. وبما أن النظام يقارن بتجزئة
ما يُقدَّم إليه، فالتجزئة المسروقة عديمة الفائدة — لا تصلح لانتحال جلسة. هذا يجعل
قاعدة البيانات هدفًا أقلّ جاذبية بكثير.
\end{secnote}

% ====================================================================
\subsection{الإبطال والانتهاء: إزاي أبوّظها بعد مرحلة معيّنة}

هذا لبّ الفصل. سؤال متكرّر: «طيّب لو عايز أبطّل جلسة في نصّها، أعمل إيه؟». الجواب
يختلف بين الرمزين لأن طبيعتهما مختلفة تمامًا.

رمز الوصول ينتهي ذاتيًّا. داخله \en{claim} اسمه \code{exp} يحمل لحظة الانتهاء (بعد
15 دقيقة). عند كل طلب يفحص المتحقّق هذا الوقت دون الرجوع إلى قاعدة البيانات؛ فإذا
مضى الوقت رُفِض الرمز. الثمن هو أننا لا نستطيع سحب رمز وصول قبل انتهائه — لكن نافذته
15 دقيقة فقط، فالخطر محدود.

رمز التحديث يُبطَل فورًا. لأنه محفوظ (مُجزَّأً) في القاعدة، يمكننا وضع علامة عليه في
أي لحظة. وله ثلاثة نطاقات إبطال:

\begin{center}
\begin{tabular}{L{4.3cm} L{3.2cm} L{6.8cm}}
\toprule
{\bfseries\color{spPrimary}سبب الإبطال} & {\bfseries\color{spPrimary}النطاق} & {\bfseries\color{spPrimary}الأثر} \\
\midrule
تسجيل الخروج & الرمز الحالي & يُعلَّم \code{IsRevoked} على هذا الرمز فقط؛ بقيّة الأجهزة تبقى. \\
كشف إعادة استخدام رمز مُبطَل & كل رموز المستخدم & اشتباه سرقة $\rightarrow$ إلغاء كل الجلسات فورًا. \\
تغيير كلمة المرور & كل رموز المستخدم & كل الجلسات القديمة تسقط؛ يلزم دخول جديد. \\
انتهاء المدّة & الرمز المعنيّ & \code{ExpiresAt} في الماضي $\rightarrow$ يُرفَض دون علامة. \\
التدوير عند الاستخدام & الرمز القديم & يُبطَل القديم لحظة إصدار بديله (لا رمزان صالحان معًا). \\
\bottomrule
\end{tabular}
\end{center}

الكود الذي ينفّذ الإبطال «الفرديّ» و«لكل العائلة». لاحظ منطق كشف إعادة الاستخدام:

\begin{english}
\begin{lstlisting}[language={[Sharp]C}]
// Reuse detection: a revoked token presented again means it was already
// rotated -- a replay attack or a stale client. Revoke the whole family.
if (stored.IsRevoked)
{
    await RevokeAllForUserAsync(stored.UserId, "Refresh token reuse detected", ct);
    return null;
}

// Global revoke: mark every live token for this user as revoked.
var tokens = await db.RefreshTokens
    .Where(t => t.UserId == userId && !t.IsRevoked)
    .ToListAsync(ct);
foreach (var t in tokens)
{
    t.IsRevoked = true;
    t.RevokedAt = clock.UtcNow;
    t.RevokedReason = reason;
}
\end{lstlisting}
\end{english}

الشرح:
\begin{itemize}
\item \code{if (stored.IsRevoked)} — لو وصلنا رمزٌ مُبطَل سلفًا، فهذا يعني أنه دُوِّر
  من قبل ثم عاد للظهور: إمّا هجوم إعادة تشغيل وإمّا عميل قديم. الأأمن أن نفترض السرقة.
\item \code{RevokeAllForUserAsync(...)} — الإبطال «لكل العائلة»: يقطع كل جلسات
  المستخدم دفعةً واحدة، فحتى لو كان المهاجم يملك رمزًا صالحًا آخر يسقط أيضًا.
\item \code{t.RevokedReason = reason} — نسجّل سبب الإبطال، فيبقى أثر يمكن تدقيقه.
\end{itemize}

الرسم التالي هو آلة حالات رمز التحديث: من الإصدار، مرورًا بالتدوير، وصولًا إلى
الحالتين النهائيتين (مُبطَل أو منتهٍ):

\begin{center}\begin{LTR}
\begin{tikzpicture}[node distance=1.3cm and 2.3cm]
  \node[tbox, text width=2.6cm] (iss) {issued\\(valid)};
  \node[sbox, text width=2.9cm, right=of iss] (rot) {rotated\\(old to revoked)};
  \node[fbox, text width=2.6cm, below=of iss] (exp) {expired};
  \node[fbox, text width=2.6cm, below=of rot] (rev) {revoked\\(terminal)};
  \draw[arr] (iss) -- node[midway,above,lbl] {on use} (rot);
  \draw[arrd] (rot) to[bend right=32] node[midway,below,lbl] {fresh pair issued} (iss);
  \draw[arr] (iss) -- node[midway,left,lbl] {ExpiresAt passed} (exp);
  \draw[arr] (iss) -- node[midway,right,lbl,align=left] {logout / reuse /\\pwd change} (rev);
  \draw[arrd] (rot.south) -- node[midway,right,lbl] {old} (rev.north);
\end{tikzpicture}
\end{LTR}\end{center}

\begin{spotlight}[ركّز هنا لإبهار اللجنة: الـ\en{refresh} مُجزَّأ وقابل للإبطال فورًا]
اجمع الصورتين معًا: رمز التحديث مُجزَّأ في القاعدة (تسريبها لا يفيد المهاجم)، وفي
الوقت نفسه قابل للإبطال لحظيًّا بثلاثة نطاقات — فرديًّا (خروج من جهاز)، ولكل العائلة
(اشتباه سرقة أو تغيير كلمة مرور)، وعالميًّا. هذا يجمع بين ميزتين نادرًا ما تجتمعان:
أمان تحت التسريب، وتحكّم كامل في الجلسة. حين تُسأل «إزاي تطرد جلسة مسروقة؟» تكون
الإجابة جاهزة: علامة واحدة على صفٍّ في القاعدة تُسقط الجلسة عند أول محاولة تجديد.
\end{spotlight}

% ====================================================================
\subsection{الدخول الموحّد (\en{SSO}) عبر \en{OAuth 2.0}}

الآن الباب الثاني للدخول: \en{Google} و\en{Microsoft}. المعيار المستخدَم هو
\en{OAuth 2.0} بنمط «كود التفويض» (\en{authorization-code}). الفكرة العامّة أن
المستخدم يثبت هويته عند المزوّد لا عندنا، فيعود متصفّحه إلينا حاملًا كودًا مؤقّتًا،
ونحن نبادله في الخلفية بمعلومات المستخدم.

\begin{idea}[الفكرة: الكود قناة آمنة]
لماذا لا يعيد المزوّد التوكن مباشرةً للمتصفّح؟ لأن المتصفّح بيئة غير موثوقة (قد يقرأ
المحتوى برنامج خبيث). بدل ذلك يعيد كودًا قصير العمر لا قيمة له وحده، ثم نبادله نحن
من الخادم — حيث يكمن سرّ تطبيقنا (\en{client secret}) — للحصول على التوكن. هكذا لا
يمرّ التوكن الحسّاس عبر المتصفّح.
\end{idea}

هذا الرسم هو التسلسل الكامل بين المتصفّح والمزوّد وواجهتنا (الأسهم مرقّمة بالترتيب):

\begin{center}\begin{LTR}
\begin{tikzpicture}
  \node[actor] (B) at (0,0) {Browser};
  \node[gbox]  (P) at (5,0) {Google /\\Microsoft};
  \node[pbox]  (A) at (10,0) {our API};
  \draw[arrd] (B.south) -- (0,-7.9);
  \draw[arrd] (P.south) -- (5,-7.9);
  \draw[arrd] (A.south) -- (10,-7.9);
  \draw[arr] (0,-1.1) -- (10,-1.1) node[midway,above,lbl] {(1) start Google sign-in};
  \draw[arr] (10,-2.0) -- (0,-2.0) node[midway,above,lbl] {(2) authorize URL + state};
  \draw[arr] (0,-2.9) -- (5,-2.9) node[midway,above,lbl] {(3) authorize + state};
  \draw[arr] (5,-3.8) -- (0,-3.8) node[midway,above,lbl] {(4) code + state};
  \draw[arr] (0,-4.7) -- (10,-4.7) node[midway,above,lbl] {(5) code + state};
  \draw[arr] (10,-5.6) -- (5,-5.6) node[midway,above,lbl] {(6) exchange code (state ok)};
  \draw[arr] (5,-6.5) -- (10,-6.5) node[midway,above,lbl] {(7) provider token + userinfo};
  \draw[arr] (10,-7.4) -- (0,-7.4) node[midway,above,lbl] {(8) OUR access + refresh};
\end{tikzpicture}
\end{LTR}\end{center}

نلاحظ في الخطوة (6) أن الواجهة تتحقّق من الـ\en{state} قبل مبادلة الكود. هذا الـ
\en{state} هو حارس مكافحة \en{CSRF}: قيمة نصدرها في الخطوة (2) ويجب أن تعود كما هي
في الخطوة (5). ندير هذا عبر \code{DataProtectionSsoStateStore}:

\begin{english}
\begin{lstlisting}[language={[Sharp]C}]
public string Issue()
{
    // payload = <expiry-unix-seconds>:<random-nonce>, protected (tamper-proof).
    var expiresAt = DateTimeOffset.UtcNow.Add(Ttl).ToUnixTimeSeconds();
    var payload   = $"{expiresAt}:{Guid.NewGuid():N}";
    return _protector.Protect(payload);          // signed + encrypted
}

public bool Validate(string? state)
{
    var payload = _protector.Unprotect(state);   // throws if forged/tampered
    var sep = payload.IndexOf(':');
    return sep > 0 && long.TryParse(payload.AsSpan(0, sep), out var exp)
        && DateTimeOffset.FromUnixTimeSeconds(exp) > DateTimeOffset.UtcNow;
}
\end{lstlisting}
\end{english}

الشرح:
\begin{itemize}
\item \code{Protect(payload)} — الـ\en{state} ليس رقمًا نحفظه في الخادم، بل حزمة
  موقّعة ومُشفّرة تحمل تاريخ انتهائها بنفسها. لا نحتاج تخزينًا من طرف الخادم، فيبقى
  التحقّق صالحًا حتى مع إعادة تشغيل التطبيق أو تعدّد نُسخه.
\item \code{Unprotect(state)} — يفشل (يرمي استثناءً) إن كان الـ\en{state} مزوّرًا
  أو معدَّلًا؛ فأيّ محاولة تزييف تسقط فورًا.
\item فحص الوقت — نرفض أي \en{state} انتهت مدّته (30 دقيقة)، فلا يُعاد استعماله لاحقًا.
\end{itemize}

بعد نجاح المبادلة وجلب بيانات المستخدم (\en{userinfo})، يتولّى
\code{SsoLoginCommandHandler} إيجاد المستخدم المحليّ أو إنشاءه، ثم — وهذه النقطة
الأهمّ — يصدر توكناتنا نحن:

\begin{english}
\begin{lstlisting}[language={[Sharp]C}]
// Resolve by the recorded external identity first, then by verified email.
var linkedUserId = await userAdministration
    .FindUserIdByExternalLoginAsync(info.Provider, info.ProviderUserId, ct);
var user = linkedUserId is { } lid
    ? await db.Users.FirstOrDefaultAsync(u => u.Id == lid, ct)
    : await db.Users.FirstOrDefaultAsync(u => u.NormalizedEmail == normalizedEmail, ct);

if (user is null) { /* first SSO sign-in: provision an Unassigned account */ }

// Link the external identity (idempotent), then issue OUR tokens.
await userAdministration.AddExternalLoginAsync(user.Id, info.Provider, info.ProviderUserId, ct);
var tokens = tokenService.IssueTokens(user, roles, user.ActiveRole, rememberMe: false);
\end{lstlisting}
\end{english}

الشرح:
\begin{itemize}
\item \code{FindUserIdByExternalLoginAsync} — نبحث أولًا بالهوية الخارجية المسجّلة
  (المزوّد $+$ معرّف المستخدم عنده)، لا بالبريد؛ فربط الحساب يصبح صريحًا ومستقرًّا
  حتى لو غيّر المزوّد البريد لاحقًا.
\item \code{if (user is null)} — أوّل دخول عبر \en{SSO} يُنشئ حسابًا محليًّا جديدًا
  (بحالة \code{Unassigned}) لأن البريد موثَّق من المزوّد.
\item \code{AddExternalLoginAsync} — نسجّل رابط الهوية الخارجية (عملية عديمة الأثر
  عند التكرار)، فيُحَل الدخول التالي إلى نفس الحساب.
\item \code{tokenService.IssueTokens(...)} — السطر الحاسم: بعد إثبات الهوية نصدر
  زوجنا نحن (\en{access} $+$ \en{refresh})، ولا نستعمل توكن \en{Google} بعد ذلك أبدًا.
\end{itemize}

\begin{spotlight}[ركّز هنا لإبهار اللجنة: بعد \en{SSO} نُصدر توكناتنا نحن]
توكن \en{Google} أو \en{Microsoft} يُستعمل مرّة واحدة فقط: لإثبات أن هذا البريد
يخصّ صاحبه. بمجرّد أن نتأكّد، نرميه ونصدر توكناتنا الخاصّة. لماذا هذا مهمّ؟ لأنه
يبقي التحكّم كاملًا في أيدينا: دورة حياة الجلسة، والإبطال بنطاقاته الثلاثة،
والأدوار — كلها تعمل بنفس المسار سواء دخل المستخدم بكلمة مرور أو عبر \en{Google}.
لا نعتمد على جلسة طرف ثالث ولا على انتهائها؛ نحن من يقرّر متى تنتهي الجلسة ومتى
تُطرَد.
\end{spotlight}

\begin{qa}[سؤال متوقّع: لو حساب جوجل اتسرق؟]
«لو حساب \en{Google} للمستخدم اتسرق، هل هيدخل عندنا؟» — الدخول عبر \en{Google} يثبت
فقط ملكيّة البريد لحظتها. لو فقد المستخدم حساب \en{Google} فالمشكلة عند \en{Google}.
أمّا عندنا فبمجرّد أن يغيّر كلمة مروره المحلية أو نُبطِل جلساته، تسقط كل التوكنات
التي أصدرناها فورًا — لأننا لا نعتمد على جلسة \en{Google}، بل على رموزنا القابلة
للإبطال.
\end{qa}

% ====================================================================
\begin{summary}[الخلاصة في سطور]
\begin{itemize}
\item \code{ApplicationDbContext} يرث \code{IdentityDbContext<ApplicationUser,
  ApplicationRole, Guid>}، فيحصل على جداول الهوية السبعة وعلاقاتها جاهزة عبر
  \code{base.OnModelCreating}؛ الوراثة من \code{DbContext} كانت ستتطلّب تعريفها كلّها يدويًّا.
\item الأدوار الخمسة (\en{SuperAdmin/Admin/Student/Consultant/ScholarshipProvider})
  تُحمَل داخل الـ\en{JWT} كـ\en{claims}، والتفويض عبر \code{[Authorize(Roles=...)]}
  مضبوط على \code{active\_role} — الدور الفعّال للجلسة.
\item رمز التحديث لا يُخزَّن خامًا، بل تُخزَّن تجزئته \en{SHA-256} في
  \code{RefreshToken.TokenHash}؛ تسريب القاعدة لا يعطي رمزًا صالحًا.
\item رمز الوصول ينتهي ذاتيًّا عبر \code{exp} (15 دقيقة، بلا رجوع للقاعدة)؛ ورمز
  التحديث يُدوَّر عند كل استخدام ويُبطَل فرديًّا/للعائلة/عالميًّا عند الخروج أو كشف
  إعادة الاستخدام أو تغيير كلمة المرور.
\item \en{SSO} يتبع \en{OAuth 2.0} (كود التفويض): كود $\rightarrow$ مبادلة
  $\rightarrow$ \en{userinfo} $\rightarrow$ إيجاد/إنشاء مستخدم، مع \code{state}
  موقّع (\code{DataProtectionSsoStateStore}) لمكافحة \en{CSRF}.
\item بعد \en{SSO} نصدر توكناتنا نحن ولا نعتمد على توكن المزوّد — تحكّم كامل في
  الجلسة والإبطال.
\end{itemize}
\end{summary}
